% R587: EGA I, 0, 3.2.1.
% Traducción: OpenAI Codex, GPT-6.1 Sol, esfuerzo Ultra.
\subsection{Prehaces sobre una base de abiertos}
\label{subsection:0.3.2-fr}

\begin{env}[3.2.1]
\label{0.3.2.1-fr}
En lo que sigue nos restringiremos a categorías $\boldsymbol{K}$
que admitan \emph{límites proyectivos} (generalizados, es decir,
correspondientes a conjuntos preordenados no necesariamente filtrantes,
cf. (T, 1.8)). Sean $X$ un espacio topológico y $\mathfrak{B}$ una base
de abiertos para la topología de $X$. Llamaremos \emph{prehaz sobre
$\mathfrak{B}$, con valores en $\boldsymbol{K}$}, a una familia de objetos
$\mathcal{F}(U)\in\boldsymbol{K}$, asociados a cada $U\in\mathfrak{B}$, y a una
familia de morfismos
$\rho_U^V:\mathcal{F}(V)\to\mathcal{F}(U)$ definidos para todo par $(U,V)$
de elementos de $\mathfrak{B}$ tales que $U\subset V$,

\oldpage[0\textsubscript{I}]{26}

con las condiciones $\rho_U^U=\text{identidad}$ y
$\rho_U^W=\rho_U^V\circ\rho_V^W$ si $U$, $V$, $W$ en $\mathfrak{B}$ son
tales que $U\subset V\subset W$. Se le puede asociar un \emph{prehaz con
valores en $\boldsymbol{K}$}~: $U\to\mathcal{F}'(U)$ en el sentido usual, tomando
para todo abierto $U$,
\[
  \mathcal{F}'(U)=\varprojlim\mathcal{F}(V),
\]
donde $V$ recorre el conjunto ordenado (para $\subset$, \emph{no filtrante} en
general) de los conjuntos $V\in\mathfrak{B}$ tales que $V\subset U$, pues los
$\mathcal{F}(V)$ forman un sistema proyectivo para los $\rho_V^W$
($V\subset W\subset U$, $V\in\mathfrak{B}$, $W\in\mathfrak{B}$). En efecto,
si $U$, $U'$ son dos abiertos de $X$ tales que $U\subset U'$, se define
$\rho'_U{}^{U'}$ como el límite proyectivo (para $V\subset U$) de los morfismos
canónicos $\mathcal{F}'(U')\to\mathcal{F}(V)$, es decir, como el único
morfismo $\mathcal{F}'(U')\to\mathcal{F}'(U)$ que, compuesto con los
morfismos canónicos $\mathcal{F}'(U)\to\mathcal{F}(V)$, da los morfismos
canónicos $\mathcal{F}'(U')\to\mathcal{F}(V)$; la comprobación de la
transitividad de los $\rho'_U{}^{U'}$ es entonces inmediata. Además, si
$U\in\mathfrak{B}$, el morfismo canónico
$\mathcal{F}'(U)\to\mathcal{F}(U)$ es un isomorfismo que permite identificar
estos dos objetos%
\footnote{Si $X$ es un espacio \emph{noetheriano}, se puede aún definir
$\mathcal{F}'(U)$ y demostrar que es un prehaz (en el sentido usual)
suponiendo solo que $\boldsymbol{K}$ admite límites proyectivos
para los sistemas proyectivos \emph{finitos}. En efecto, si $U$ es un abierto
cualquiera de $X$, existe un recubrimiento \emph{finito} $(V_i)$ de $U$ formado
por conjuntos de $\mathfrak{B}$; para todo par $(i,j)$ de índices, sea
$(V_{ijk})$ un recubrimiento finito de $V_i\cap V_j$ formado por conjuntos de
$\mathfrak{B}$. Sea $I$ el conjunto formado por los $i$ y las ternas $(i,j,k)$,
ordenado por las únicas relaciones $i>(i,j,k)$, $j>(i,j,k)$; se toma entonces
como $\mathcal{F}'(U)$ el límite proyectivo del sistema de los $\mathcal{F}(V_i)$
y $\mathcal{F}(V_{ijk})$; se comprueba fácilmente que no depende de los
recubrimientos $(V_i)$ y $(V_{ijk})$ y que $U\to\mathcal{F}'(U)$ es un
prehaz.}.
\end{env}
